% quadprog.tex — companion paper for cvx-quadprog
\documentclass[11pt,a4paper]{article}

% --------------------------------------------------------------------------
% Packages
% --------------------------------------------------------------------------
\usepackage[T1]{fontenc}
\usepackage[utf8]{inputenc}
\usepackage{lmodern}
\usepackage{microtype}
\usepackage[a4paper, margin=2.5cm]{geometry}

\usepackage{amsmath}
\usepackage{amssymb}
\usepackage{amsthm}

\usepackage{booktabs}
\usepackage{array}

\usepackage{xcolor}
\usepackage{enumitem}

\usepackage[round]{natbib}
\usepackage[colorlinks=true, linkcolor=blue!60!black, citecolor=blue!60!black,
            urlcolor=blue!60!black]{hyperref}

% --------------------------------------------------------------------------
% Styling
% --------------------------------------------------------------------------
\theoremstyle{plain}
\newtheorem{proposition}{Proposition}
\theoremstyle{remark}
\newtheorem*{remark}{Remark}

\newcommand{\R}{\mathbb{R}}
\newcommand{\T}{^{\mathsf{T}}}
\DeclareMathOperator*{\argmin}{arg\,min}

\title{\bf A Dual Active-Set Quadratic Programming Solver\\ in NumPy and SciPy}

\author{%
  Thomas Schmelzer\\[2pt]
  \small Jebel Quant Research\\
  \small \texttt{https://github.com/Jebel-Quant/quadprog}
}

\date{August 2026}

\begin{document}
\maketitle

\begin{abstract}
\noindent
The Goldfarb--Idnani dual active-set method \citep{goldfarb1983} remains one of
the most effective algorithms for small and medium dense strictly convex
quadratic programs, and the \texttt{quadprog} lineage descending from Powell's
\textsc{zqpcvx} \citep{powell1983} and Turlach's Fortran translation
\citep{turlach2019} is still the reference implementation in both R and Python.
That lineage is compiled code: it presumes a C toolchain and a build step. We
describe a reimplementation in NumPy and SciPy alone, and use it as a controlled
setting in which to ask which of the reference's implementation decisions are
essential to the algorithm and which are artefacts of the language it was written
in.

Three findings are reported. First, the reference's chain of Givens rotations for
adding a constraint can be replaced by a single Householder reflection without
altering the sequence of iterates, because the quantities the solver consumes are
invariant to the choice of orthogonal reduction; we prove this and confirm it
empirically on 3027 problems, where both implementations agree on every iteration
count. Second, the reference's packed storage for the triangular factor, easily
mistaken for a memory optimisation, is load-bearing for performance: it is what
keeps the active submatrix contiguous and therefore admissible to a BLAS packed
solve, worth a factor of 10.2 on that operation. Third, detecting that a
constraint column holds a single nonzero --- which is what a bound constraint is
--- reduces three per-iteration products from $O(n^2)$ or $O(nm)$ to indexing.

The result is faster than the compiled reference for $n \gtrsim 135$, by
11 times at $n = 1600$, while remaining slower for small problems by a
factor set by interpreter dispatch rather than by arithmetic. Attacking that
floor at its source --- the \emph{number} of iterations rather than their cost
--- by guessing the active set and certifying it against the KKT conditions moves
the crossover to $n \approx 65$ and the margin to 68 times at $n = 1600$.
\end{abstract}

\section{Introduction}
\label{sec:intro}

The strictly convex quadratic program
\begin{equation}
  \min_{x \in \R^n} \; \tfrac12 x\T G x - a\T x
  \qquad \text{subject to} \qquad C\T x \ge b,
  \label{eq:qp}
\end{equation}
with $G \in \R^{n \times n}$ symmetric positive definite,
$C \in \R^{n \times m}$ and the first $m_{\mathrm{eq}}$ constraints understood as
equalities, is among the most thoroughly studied problems in numerical
optimisation \citep{nocedal2006,fletcher1987,boyd2004}. It is also among the most
frequently solved in practice: mean--variance portfolio selection
\citep{markowitz1952} is exactly~\eqref{eq:qp}, and so is every step of a
sequential quadratic programming method and every horizon of a linear model
predictive controller.

For large sparse instances the field has moved decisively towards operator
splitting \citep{stellato2020} and interior-point methods. For the dense
small-to-medium regime, however --- say $n$ from a handful to a few thousand,
with $m$ of the same order --- active-set methods retain two decisive advantages:
they terminate at an exactly feasible vertex of the active-set lattice rather
than approaching it asymptotically, and they warm-start almost perfectly, which
is why parametric solvers such as qpOASES \citep{ferreau2014} are built on them.

Within that regime the Goldfarb--Idnani (GI) dual method
\citep{goldfarb1982,goldfarb1983} is the natural choice. Its distinguishing
property is that it maintains \emph{dual} feasibility from the first iterate and
drives primal infeasibility to zero, rather than the reverse. Because the
unconstrained minimiser $G^{-1}a$ is trivially dual feasible, the method needs no
phase-one problem at all --- a considerable practical simplification, and the
reason GI is preferred over primal active-set methods for this problem class.

\subsection{The reference implementation and its lineage}

The canonical implementation has an unusually well-documented pedigree. Powell
distributed \textsc{zqpcvx} \citep{powell1983} shortly after the paper appeared;
Turlach translated and repackaged the method as the Fortran core of the R package
\texttt{quadprog} \citep{turlach2019}, which remains the standard R interface;
that Fortran was subsequently transliterated to C and wrapped with Cython
\citep{cython2011} as the Python package \texttt{quadprog}
\citep{quadprogpy}. The numerical core is recognisably the same code in all
three, down to variable names and the comment quoting Powell on the machine
epsilon.

That code is good code, and it is fast. But it is compiled code, and it carries
the consequences: a source install needs a C compiler, wheels must be built for
every platform and Python version, and the implementation is opaque to the
majority of its users, who are Python programmers rather than C programmers.

\subsection{Contribution}

We reimplement GI using only NumPy \citep{numpy2020} and SciPy
\citep{scipy2020}. The motivation is partly practical --- no compiler, no build
step, and a solver that can be read and modified by the people who use it --- but
the reimplementation also creates a controlled experiment. The reference makes a
series of implementation choices which are natural in Fortran or C and which have
been carried forward for four decades without, as far as we are aware, being
revisited. Reimplementing in a language with completely different cost structure
forces each of them to be re-derived, and it turns out that some are essential to
the algorithm, some are essential for reasons other than the ones usually given,
and some are artefacts.

Concretely, this paper contributes:
\begin{itemize}[leftmargin=1.4em,itemsep=2pt]
  \item a proof that the constraint-insertion update may use any orthogonal
        reduction --- Householder or Givens --- without changing the sequence of
        iterates the solver visits (\S\ref{sec:invariance}), together with the
        empirical confirmation that iteration counts match the reference exactly
        on every problem tested;
  \item the observation that the reference's packed triangular storage is a
        performance feature rather than a memory feature, because it is what
        makes the active submatrix contiguous and hence eligible for a BLAS
        packed triangular solve (\S\ref{sec:packed});
  \item a per-column structure detection that reduces three per-iteration
        products to indexing when a constraint is a bound (\S\ref{sec:structure});
  \item a differential validation methodology which compares not merely
        minimisers but complete active-set trajectories against the reference
        (\S\ref{sec:agreement}).
\end{itemize}

\section{The dual active-set method}
\label{sec:method}

\subsection{Optimality conditions}

Let $\mathcal{A} \subseteq \{1,\dots,m\}$ index a working set of constraints and
write $A = C_{\cdot\mathcal{A}}$ for the matrix of their normals. The
Karush--Kuhn--Tucker conditions for~\eqref{eq:qp} require a primal--dual pair
$(x, u)$ with
\begin{align}
  G x - a &= C u, \label{eq:stat}\\
  C\T x - b &\ge 0, \qquad u_i \ge 0 \ \ (i > m_{\mathrm{eq}}), \label{eq:feas}\\
  u_i \, (C\T x - b)_i &= 0 \quad \text{for all } i. \label{eq:comp}
\end{align}
A primal active-set method maintains~\eqref{eq:feas} and works towards
\eqref{eq:stat}; the dual method maintains \eqref{eq:stat},
\eqref{eq:comp} and the sign conditions, and works towards primal feasibility.

\subsection{One iteration}

GI keeps an iterate $x$ that minimises the objective subject to the working-set
constraints holding with equality, together with the corresponding multipliers
$u$. If every constraint is satisfied at $x$, the KKT conditions hold and the
method stops. Otherwise it selects the most violated constraint --- scaled by the
norm of its normal, so that the choice is invariant to how each constraint
happens to be written --- and moves $x$ towards that constraint's boundary.

Two step lengths compete. The primal step $t_2$ brings the violated constraint's
slack to zero. The dual step $t_1$ is the largest step for which no multiplier of
an active inequality goes negative. If $t_1 < t_2$ the method takes a partial
step, drops the constraint whose multiplier reached zero, and repeats; otherwise
it takes the full step and adds the new constraint to the working set. Since the
objective increases monotonically and each iteration adds a constraint, the
method terminates finitely.

The linear algebra per iteration is where all the cost lies. Writing
$G^{-1} = J J\T$ with $J$ upper triangular --- so $J = R_G^{-1}$ for the Cholesky
factor $G = R_G\T R_G$ --- and letting $n_\ast$ denote the normal of the entering
constraint, the method needs
\begin{align}
  d &= J\T n_\ast, \label{eq:d}\\
  z &= J_2 d_2, \label{eq:z}\\
  r &= R^{-1} d_1, \label{eq:r}
\end{align}
where $J = [\,J_1 \; J_2\,]$ and $d = (d_1, d_2)$ are split at the size $k$ of
the working set, and $R$ is the upper triangular factor of the QR factorisation
of $J\T A$. Here $z$ is the primal search direction --- the component of the
constraint normal orthogonal, in the $G^{-1}$ metric, to the active constraints
--- and $-r$ is the corresponding direction for the multipliers.

The essential structural fact, and the reason GI is efficient, is that
\begin{equation}
  J\T A = \begin{bmatrix} R \\ 0 \end{bmatrix}
  \label{eq:invariant}
\end{equation}
can be \emph{updated} rather than recomputed when a column enters or leaves $A$.
Updating costs $O(n^2)$; recomputing costs $O(nk^2)$. Over a run of $O(n)$
iterations this is the difference between $O(n^3)$ and $O(n^4)$, which is why the
factor is carried and updated rather than rebuilt.

\begin{remark}
Equation~\eqref{eq:r} is a least-squares solve in disguise. Since
$A\T G^{-1} A = (J\T A)\T (J\T A) = R\T R$ and $A\T G^{-1} n_\ast = R\T d_1$,
\begin{equation}
  r = R^{-1} d_1 = (R\T R)^{-1} R\T d_1
    = (A\T G^{-1} A)^{-1} A\T G^{-1} n_\ast
    = \argmin_y \; \lVert (J\T A) y - J\T n_\ast \rVert_2 ,
  \label{eq:lsq}
\end{equation}
and $z$ is its residual mapped back through $J$. It is tempting to conclude that
a library least-squares routine should simply be called here. It should not: the
quantities it would recompute are already available in factored form, and forming
them afresh at every iteration is far more expensive than updating them.
\end{remark}

\section{Reimplementation in an interpreted setting}
\label{sec:reimpl}

\subsection{What changes}

The reference is written for a machine model in which a scalar operation is
cheap and the unit of composition is the loop. NumPy inverts this: a scalar
operation reached through the interpreter costs on the order of a microsecond,
while an operation on a whole array costs almost the same. Every loop over
elements must therefore become an array expression, and the design question
becomes not ``how many floating-point operations?'' but ``how many array
operations, and does each reach a tuned kernel?''

Two consequences follow immediately. First, the reference's hand-written
\texttt{dot} and \texttt{axpy} --- a private BLAS level-1 \citep{blas1979} in
\texttt{linear-algebra.c} --- should be replaced not by NumPy equivalents but by
calls that reach the vendor BLAS level-2 \citep{blas2level} and LAPACK
\citep{lapack1999}. Second, any part of the algorithm that is irreducibly
sequential in the number of \emph{iterations} rather than in the number of
elements will remain interpreted, and will dominate for small $n$. Both
predictions are borne out (\S\ref{sec:perf}).

\subsection{Design constraints}

We adopt three constraints. The public signature is kept identical to
\citep{quadprogpy}, including the argument order, the convention that the linear
term is subtracted, the column-wise $\ge$ form of the constraints, and the text
of the error messages, so that the package is a drop-in replacement. Input arrays
are never modified, unlike the reference which destroys $G$ and $a$. And the
dependency set is exactly NumPy and SciPy: no compiler, no Cython, no build step.

\section{Householder versus Givens for the insertion}
\label{sec:invariance}

\subsection{Two reductions}

When a constraint enters the working set, the vector $d$ of~\eqref{eq:d} must be
reduced so that its components beyond the $k$-th vanish, and the same orthogonal
transformation must be applied to the columns of $J$ so that
\eqref{eq:invariant} is preserved.

The reference does this with a chain of Givens rotations \citep{givens1958}: one
rotation per trailing component, each mixing two adjacent columns of $J$ over all
$n$ rows. In C this is excellent --- the inner loop is tight, the rotations are
applied in place, and the arithmetic is minimal. In an interpreted setting it is
close to the worst possible structure, because it is $O(n)$ interpreter
round-trips per insertion, each doing only $O(n)$ arithmetic. Profiling our first
faithful transliteration put 85\% of runtime at $n = 150$ in exactly
this loop.

The alternative is a single Householder reflection \citep{householder1958},
\begin{equation}
  W = I - \frac{2}{\lVert v \rVert^2} v v\T ,
  \qquad v = d_{k:} - \alpha e_1, \qquad \alpha = \pm \lVert d_{k:} \rVert ,
\end{equation}
which achieves the same reduction in one matrix--vector product and one rank-one
update, independently of $n - k$. Both reach the BLAS: \texttt{gemv} and
\texttt{ger} respectively, the latter writing directly into $J$'s buffer so that
no $O(n(n-k))$ temporary is allocated. We take the sign of $\alpha$ from $d_k$,
matching the reference's convention, and form the leading entry of $v$ by the
cancellation-free rearrangement
\begin{equation}
  v_1 = d_k - \operatorname{sign}(d_k)\lVert d_{k:} \rVert
      = \frac{-\operatorname{sign}(d_k)\,\lVert d_{k+1:} \rVert^2}
             {|d_k| + \lVert d_{k:} \rVert},
\end{equation}
which avoids the subtractive cancellation that the direct formula suffers when
$d_{k:}$ is already nearly a positive multiple of $e_1$ \citep{golub2013}.

Deletion keeps the Givens chase, and the reason no analogous single-transformation
form exists is worth making precise. Removing a column shifts the remainder left
and leaves a single subdiagonal in the trailing block. A Householder reflection
annihilates all but one component of a \emph{single} vector; those subdiagonal
entries lie in distinct columns, so no one reflection reaches them. The chase is
additionally sequential --- each rotation's parameters are read from entries the
previous rotation wrote --- which rules out computing them in a batch.

What the chase does admit is the same routing to the BLAS. Each step applies a
$2 \times 2$ transformation to a pair of adjacent columns of $J$; spelled out in
NumPy this allocates two $n$-vectors per step, while BLAS \texttt{rot} performs it
in place. Because $J$ is stored Fortran-ordered the two columns are contiguous,
so the call writes through to $J$'s own buffer. The pairwise update is 2.0 to
3.0 times faster for $n$ between 50 and 2000, and a complete solve is 5 to
8\% faster on problems whose working set churns; where no constraint is ever
dropped the change is inert. We note that this is a smaller effect than either
\S\ref{sec:packed} or \S\ref{sec:structure}, and that we found it only after a
first measurement showed no gain at all --- on a test set in which, as profiling
established, no deletion ever occurred.

A detail here corroborates Proposition~\ref{prop:invariance} a second time, on the
deletion side. BLAS offers no $2 \times 2$ reflection, only the rotation
$\left[\begin{smallmatrix} c & -s \\ s & c \end{smallmatrix}\right]$, whose second
output is the exact negation of the reflection's; one sign flip therefore recovers
the reference's convention. Proposition~\ref{prop:invariance} says that flip is
unnecessary, since using the rotation outright replaces $R$ by $S_k R$ and $J$ by
$J S$. We implemented that variant and confirmed it: every iteration count in
\S\ref{sec:agreement} is reproduced, with $|R|$ and $|J|$ elementwise unchanged
and only signs differing. It is not retained. It measures no faster --- the flip
is one in-place pass --- and it forfeits the symmetry of the reflection, which is
what allows one $2 \times 2$ to serve both $J$'s columns and $R$'s rows; the
rotation requires its transpose on the $R$ side and turns the degenerate case into
a signed exchange.

\subsection{The transformations differ; the iterates do not}

Householder and Givens reductions do not produce the same $Q$ and $R$: signs
along the diagonal of $R$, and hence the signs of some columns of $J$, differ.
Since the solver consumes $R$ and $J$, one might expect the substitution to
perturb the computation. It does not, and the reason is worth stating precisely.

\begin{proposition}[Invariance of the iterates]
\label{prop:invariance}
Let $\hat{J} = J S$ and $\hat{R} = S_k R$ where $S = \operatorname{diag}(s)$ with
$s_i \in \{-1, +1\}$ and $S_k$ is its leading $k \times k$ block. Then the
quantities $r$ and $z$ of \eqref{eq:r} and \eqref{eq:z} are unchanged, and in
IEEE arithmetic they are unchanged \emph{exactly}.
\end{proposition}

\begin{proof}
From \eqref{eq:lsq}, $r = (A\T G^{-1} A)^{-1} A\T G^{-1} n_\ast$ depends only on
$A$, $n_\ast$ and $G$, and none of these is altered by the choice of reduction.
Directly: $\hat{d} = \hat{J}\T n_\ast = S d$, so
$\hat{r} = \hat{R}^{-1}\hat{d}_1 = (S_k R)^{-1} S_k d_1 = R^{-1} d_1 = r$. For
the primal direction,
$\hat{z} = \hat{J}_2 \hat{d}_2 = \sum_{j > k} (s_j J_{\cdot j})(s_j d_j)
= \sum_{j>k} s_j^2 J_{\cdot j} d_j = z$. Both cancellations are multiplications
by $\pm 1$, which are exact in IEEE~754, so the identities hold in floating point
and not merely in exact arithmetic.
\end{proof}

The proposition explains why the substitution is safe, but it does not by itself
guarantee that the two implementations follow the same trajectory: rounding in
the reductions themselves still differs, and the algorithm's branches --- which
constraint is most violated, whether $t_1 \ge t_2$ --- are discontinuous
functions of those roundings. That the trajectories coincide anyway is an
empirical claim, and we test it directly in \S\ref{sec:agreement} by comparing
iteration counts rather than only solutions.

\begin{remark}
Proposition~\ref{prop:invariance} also disposes of a natural objection to the
whole enterprise. One might ask why the solver maintains $J$, an $n \times n$
dense matrix, rather than the orthogonal factor $Q$ of $J\T A$ together with the
triangular $R_G$. The answer is that $J$ folds $R_G^{-1}$ and the accumulated $Q$
into a single matrix, so that \eqref{eq:d} and \eqref{eq:z} are each one
\texttt{gemv} regardless of $k$.
\end{remark}

\section{Storage as a performance decision}
\label{sec:packed}

The reference stores $R$ as packed columns: entry $(i,j)$, $i \le j$, at flat
offset $j(j+1)/2 + i$. The obvious reading is that this halves memory, and our
first reimplementation accordingly replaced it with a dense $(r, r)$ array,
documenting the change as affecting ``only the addressing''. That was wrong, and
the way in which it was wrong is instructive.

Equation~\eqref{eq:r} is solved once per iteration against the leading $k \times
k$ block of $R$. In a dense array that block is a \emph{strided} view: its
column stride is $r$, not $k$. Handed such a view, SciPy's LAPACK wrapper cannot
pass it through and copies it, at a cost of $O(k^2)$ memory traffic per
iteration. In packed storage the same block is the leading $k(k+1)/2$ entries ---
contiguous by construction, at every $k$ --- and BLAS \texttt{tpsv} reads it in
place.

The effect is out of all proportion to the arithmetic. At $n = 700$ with
$k = 383$:

\begin{center}
\begin{tabular}{lr}
  \toprule
  Triangular solve variant & Time ($\mu$s) \\
  \midrule
  \texttt{trtrs} on the strided view (dense $R$) & 77.0 \\
  \texttt{trtrs} on a contiguous copy            & 22.6 \\
  \texttt{tpsv} on the packed triangle           & 7.5 \\
  \bottomrule
\end{tabular}
\end{center}

\noindent
A factor of 10.2, none of it arithmetic. The giveaway in the profile was
that this solve was running at roughly 1.2\,Gflop/s while a neighbouring
\texttt{gemv} on the same data managed 10\,Gflop/s. Over a complete
solve the operation fell from 15.5\,ms
(21\% of runtime) to 2.0\,ms (3.5\%).

The cost is borne by the deletion routine, which mixes two \emph{rows} of $R$
across a range of columns. Because column offsets grow with the column index,
what is a contiguous slice in the dense layout becomes a gather in the packed
one. Insertion is unaffected: a column is one contiguous run in either layout.

We draw a general moral. In compiled numerical code, data layout is chosen for
cache behaviour and for memory. In array-language numerical code it is
additionally, and often dominantly, a question of \emph{which kernels the layout
makes reachable}. A layout that forces a copy at the library boundary can cost an
order of magnitude on an operation whose flop count is negligible.

\section{Exploiting constraint structure}
\label{sec:structure}

A bound constraint $x_i \ge \ell_i$ is a single nonzero in its column of $C$, and
a box-constrained problem consists of nothing else. When the entering
constraint's column is $\nu e_\rho$ for a scalar $\nu$ and index $\rho$, three of
the per-iteration products degenerate:

\begin{center}
\begin{tabular}{lll}
  \toprule
  Quantity & General & Column is $\nu e_\rho$ \\
  \midrule
  slack $C\T x$          & $O(nm)$ \texttt{gemv} & $O(m)$ gather \\
  $d = J\T n_\ast$       & $O(n^2)$ \texttt{gemv} & $O(n)$: one scaled row of $J$ \\
  $z\T n_\ast$           & $O(n)$ dot            & $O(1)$ \\
  \bottomrule
\end{tabular}
\end{center}

The slack product matters most, because a box-constrained problem has $m = 2n$
and the product is evaluated once per outer iteration; before this change it was
the single largest consumer of arithmetic in the solver.

Detection is performed once, per column. The per-column choice is not a
refinement but the point: the case worth optimising is \emph{mixed}. A
mean--variance problem \citep{markowitz1952} carries a dense budget column
$\sum_i x_i = 1$ alongside $2n$ bounds, and an all-or-nothing test would observe
that single dense column and send the entire problem down the general path. The
slack evaluation is additionally given a three-way choice --- an all-unit gather,
a compiled sparse CSR product, or the dense product --- since a mixed problem of
this shape has density $O(1/n)$ and is far too sparse to want the last.

Because this is a fast path around arithmetic the general path performs anyway,
it cannot alter the answer; we verify that claim rather than assert it
(\S\ref{sec:agreement}). One trap deserves naming: an all-zero column must not be
mistaken for a unit column. It expresses $0 \ge b_i$, which no $x$ can influence,
and treating it as a unit column would scale a row of $J$ by zero and lose the
infeasibility verdict.

\section{Experiments}
\label{sec:experiments}

\subsection{Methodology}

All measurements are on an \texttt{arm64} machine, CPython 3.12, NumPy 2.5.1,
SciPy 1.18, against \texttt{quadprog} 0.1.13 built from its published source.
Timings are the best of five batches after a warm-up call; we report the best
rather than the mean because the quantity of interest is the achievable time and
the noise is one-sided.

Random problems are generated with $G = BB\T + nI$ for $B$ Gaussian, which is
positive definite by construction and moderately conditioned. Where box
constraints are used, the bounds straddle the unconstrained minimiser so that
roughly half of them bind at the solution.

\subsection{Agreement with the reference}
\label{sec:agreement}

Correctness for a solver of this kind is usually argued by comparing minimisers.
That is too weak a test to detect the class of error we are most exposed to: a
change that perturbs the active-set trajectory while still converging to the
right answer would pass, and would be a latent hazard on degenerate problems. We
therefore compare complete trajectories, by requiring that both implementations
report identical iteration counts --- the number of constraints added and the
number dropped --- as well as identical minimisers, objectives, multipliers and
active sets.

Over 4000 random problems with $2 \le n \le 11$, up to 14 constraints and mixed
equality counts:

\begin{center}
\begin{tabular}{ll}
  \toprule
  Quantity & Agreement with \citep{quadprogpy} \\
  \midrule
  Iteration counts (added, dropped) & exact, 3027 / 3027 feasible problems \\
  Infeasibility verdict             & exact, 973 / 973 infeasible problems \\
  Minimiser $x$                     & max abs.\ difference 3.0e-09 \\
  Objective                         & max rel.\ difference 2.5e-12 \\
  \bottomrule
\end{tabular}
\end{center}

\noindent
That the iteration counts agree exactly on every feasible problem, and the
infeasibility verdict on every infeasible one, is a considerably stronger
statement than agreement on the answer: the two implementations add and drop the
same constraints in the same order, despite one using Householder reflections and
the other Givens rotations. Proposition~\ref{prop:invariance} explains why this
is possible; the experiment establishes that it in fact occurs.

The suite also includes the upstream test problems verbatim, with their published
iteration counts asserted, and a differential test at $n = 60, 140, 220$ where
the vectorised paths dominate. The implementation carries 1066 tests at
100\% line and branch coverage.

\subsection{Where the two legitimately differ}

Duplicated or linearly dependent constraints make the \emph{dual} solution
non-unique: a multiplier may sit on either copy. Both implementations return a
valid KKT point but not necessarily the same one, so multipliers and reported
active sets differ while $x$ and the objective do not. We detected this as an
apparent test failure and, rather than loosen a tolerance, verified that both
answers satisfy \eqref{eq:stat}--\eqref{eq:comp} to machine precision --- the
stationarity residual is $\sim\!1e-16$ for both --- and changed the test to
check the KKT conditions on such problems instead of demanding an identical dual.

\subsection{Accuracy}

Both implementations accumulate the objective incrementally rather than
re-evaluating the quadratic, following \citep{goldfarb1983}. Measuring each
against a direct re-evaluation at \emph{its own} minimiser over 2164 problems,
the worst-case drift is 1.5e-8 absolute (7.4e-15 relative) here
against 3.7e-8 (1.8e-14) for the reference; but neither dominates
problem by problem, the reference being closer on 801 problems and this
implementation on 782, with 581 ties.

For the KKT residual proper, over 3000 problems the worst relative stationarity
residual $\lVert Gx - a - Cu \rVert_\infty$, scaled, is 7.3e-13 here
against 8.8e-13 for the reference, with this implementation strictly more
accurate on 1035 problems and the reference on 869. The two are therefore of
equal numerical quality, as the backward stability of both reductions
\citep{wilkinson1965,higham2002} would predict.

\begin{remark}
An earlier version of the objective experiment evaluated the exact objective at
\emph{this} implementation's minimiser and compared both accumulated values
against that single point, which charged the reference for the distance between
the two minimisers in addition to its own drift and produced a flattering
1.4e-8 versus 3.1e-7. We record the error because its direction was
self-serving and it survived a code change unnoticed, being computed outside the
test suite.
\end{remark}

\subsection{Performance}
\label{sec:perf}

Box-constrained problems, $n$ variables and $2n$ constraints, per solve:

\begin{center}
\begin{tabular}{r r r l}
  \toprule
  $n$ & This work (ms) & {\texttt{quadprog} C (ms)} & Ratio \\
  \midrule
  10   & 0.077 & 0.006 & 12.5$\times$ slower \\
  25   & 0.16  & 0.017 & 9.4$\times$ slower \\
  50   & 0.40  & 0.076 & 5.3$\times$ slower \\
  100  & 0.96  & 0.60  & 1.6$\times$ slower \\
  200  & 2.8   & 5.5   & \textbf{2.0$\times$ faster} \\
  400  & 11.5  & 47    & \textbf{4.1$\times$ faster} \\
  800  & 53    & 461   & \textbf{8.8$\times$ faster} \\
  1600 & 374   & 4121  & \textbf{11$\times$ faster} \\
  \bottomrule
\end{tabular}
\end{center}

The crossover is at $n \approx 135$, located by sweeping the interval: the ratio
passes unity between $n = 130$ ($1.02\times$) and $n = 140$ ($0.92\times$).

Below the crossover the cost is interpreter dispatch: roughly 14\,$\mu$s
per iteration spread over some fourteen array operations, against
6\,$\mu$s for the reference to complete an entire $n = 10$ solve.
This is a floor set by the language implementation, not by the algorithm.

Above the crossover this implementation is the faster of the two, and the reason
is not subtle: the reference's inner products and \texttt{axpy}s are hand-written
scalar loops in \texttt{linear-algebra.c}, while the corresponding work here is
expressed as BLAS calls that reach vendor-tuned, vectorised kernels. The
asymptotic behaviour is $O(n^3)$ for both; the constant differs by an order of
magnitude in the vendor's favour.

After both optimisations, the residual profile at $n = 700$ is dominated by the
insertion update, at roughly 50\% of runtime.

\subsection{Attacking the iteration count instead}
\label{sec:pdas}

Every measurement above accepts the dual method's defining feature: it reaches the
active set one constraint at a time, so its iteration count grows with the size of
that set --- 74 outer iterations at $n = 100$ on a budget-plus-bounds problem.
Below the crossover a solve costs very nearly its count of interpreter-level
operations, so the floor identified above is a floor on \emph{iterations times
dispatch}, and the preceding sections attacked only the second factor.

The first factor yields to a different method. A primal-dual active set guesses
the whole active set $\mathcal{A}$ at once, solves the equality-constrained
problem it induces,
\[
  \bigl(C_\mathcal{A}\T G^{-1} C_\mathcal{A}\bigr)\lambda
    = b_\mathcal{A} - C_\mathcal{A}\T G^{-1} a,
  \qquad
  x = G^{-1}a + G^{-1} C_\mathcal{A} \lambda,
\]
and repairs the guess from the signs that come back: an active inequality wanting
a negative multiplier leaves the set, a violated inactive constraint joins it.
Across every family we measured this settles in two to four repairs
\emph{independently of $n$}. It trades flops, which are abundant at these sizes,
for dispatches, which are not.

The method is not globally convergent, so the trade is only sound because the
answer is checked. The program is strictly convex, so the KKT conditions are
sufficient and not merely necessary: a point satisfying them is \emph{the}
minimiser. Every candidate is therefore certified before it is returned, and one
that fails is discarded in favour of the exact walk. This is not a formality. Over
1164 converged attempts, two had settled on a set that was not optimal --- one of
them $0.85$ from the true minimiser --- and the certificate rejected both; among
the points it accepted, the largest disagreement with the exact walk was
$3.3\times10^{-15}$.

\begin{center}
\begin{tabular}{r r r l}
  \toprule
  $n$ & Certified path (ms) & {\texttt{quadprog} C (ms)} & Ratio \\
  \midrule
  10   & 0.081 & 0.006 & 13.2$\times$ slower \\
  25   & 0.11  & 0.017 & 6.4$\times$ slower \\
  50   & 0.14  & 0.076 & 1.9$\times$ slower \\
  100  & 0.24  & 0.60  & \textbf{2.6$\times$ faster} \\
  200  & 0.56  & 5.5   & \textbf{9.7$\times$ faster} \\
  400  & 2.4   & 47    & \textbf{19$\times$ faster} \\
  800  & 13.4  & 461   & \textbf{34$\times$ faster} \\
  1600 & 61    & 4121  & \textbf{68$\times$ faster} \\
  \bottomrule
\end{tabular}
\end{center}

The crossover moves from $n \approx 135$ to $n \approx 65$, the ratio passing
unity between $n = 60$ ($1.21\times$) and $n = 70$ ($0.84\times$). It lands that
early for a reason worth stating: the reference is itself a dual active-set walk,
so it too adds one constraint per iteration --- roughly $0.45n$ of them on these
problems --- where the guess converges in about three repairs whatever $n$ is.
Each repair is far heavier, but a heavy constant beats a light $O(n)$.

Below twelve variables the path is not attempted: the guess is likeliest to be
linearly dependent when there are few variables to spread the active set over, and
that is also where it wins least. Measured over 300 instances per size, the
certified fraction is 100\% from twelve variables upward on all four families,
against 23\% at $n = 3$ on equality-constrained problems. It is offered as an
option rather than a default, because the working-set edits it counts are not the
reference's iteration counts, and \S\ref{sec:agreement} is a claim about those.

\section{Discussion}
\label{sec:discussion}

Three observations generalise beyond this solver.

\paragraph{Layout determines which kernels are reachable.}
The packed-storage result (\S\ref{sec:packed}) was worth a factor of ten on its
operation and none of it was arithmetic. In compiled code, data layout is chosen
for cache behaviour; in array-language code it additionally decides whether a
library call can proceed in place or must copy at the boundary. A strided view
handed to a wrapper that requires contiguity is a silent $O(k^2)$ tax per call.
We had inherited the right layout from the reference, discarded it as a mere
memory optimisation, and had to rediscover its purpose by profiling.

\paragraph{Flop counts mislead in both directions.}
Our initial accounting identified the slack product as 44\% of
arithmetic at $n = 400$ and predicted a large gain from eliminating it; the
realised gain was 1.35 times, because the eliminated work was
vectorised and therefore cheap per flop. Conversely the triangular solve was a
negligible 0.5\% of arithmetic and 21\% of runtime. In
the implicit-$Q$ experiment the flop count was the correct guide only once it was
computed properly, at which point it predicted the observed loss. The reliable
procedure is to measure per-operation wall-clock inside the real iteration, which
is what eventually located both wins.

\paragraph{Invariance arguments license aggressive substitution.}
Proposition~\ref{prop:invariance} is elementary, but without it the Householder
substitution would have been an unjustifiable risk in a solver whose branches are
discontinuous in its own rounding. Identifying which quantities an iterative
method actually consumes --- as opposed to which it happens to compute --- is
what makes it safe to change the representation underneath. We would expect
similar arguments to license similar substitutions in other active-set and
projection methods.

\subsection{Limitations}

The implementation targets dense problems and inherits GI's assumption that $G$
is positive definite; semidefinite $G$ requires regularisation or a different
method. Small problems remain slower than the compiled reference by up to
13 times, and closing that gap appears to require leaving pure NumPy. The structure detection of \S\ref{sec:structure}
recognises single-nonzero columns and general sparsity in the slack product but
does not exploit sparsity in $G$, for which a different factorisation strategy
would be appropriate. Finally, our differential validation is against one
reference implementation; agreement on trajectories is strong evidence of
faithfulness to that implementation, not proof of correctness in the abstract.

\section{Conclusion}

The Goldfarb--Idnani method is forty years old, and its canonical
implementation is scarcely younger. Reimplementing it in an array language proved to be a useful
way of separating the algorithm from its Fortran-era encoding. The chain of
Givens rotations turned out to be replaceable, and provably so. The packed
triangular storage turned out to be indispensable, though not for the reason it
is usually given. The one structural fact the reference does not exploit --- that
a bound constraint is a single nonzero --- was worth another 25\%.
Two further attractive ideas turned out to be worth nothing at all, and we have
reported them at the same length as the successes.

The resulting solver needs no compiler and is, for $n \gtrsim 135$, several times
faster than the compiled code it reimplements.

\section*{Reproducibility}

The implementation, the complete test suite, and the experiment scripts
underlying every table in this paper are available at
\url{https://github.com/Jebel-Quant/quadprog}. The differential tests against
\citep{quadprogpy} run in continuous integration on CPython 3.11, 3.12, 3.13 and
3.14, on Linux, macOS and Windows.

\bibliographystyle{plainnat}
\bibliography{quadprog}

\end{document}
